Achieving sustained recursive self-improvement requires more than increasingly autonomous research agents, while it also requires an infrastructure through which research experience can be reliably recorded, verified, inherited, and reused across generations. Agent-Native Research Lab focuses on this infrastructural layer, summarized in Figure~\ref{fig:bt_amber}. 

$\bullet$ \textbf{(1) Human-Centric vs. Agent-Native Research Artifact.} Conventional scientific papers are designed for human communication and would discard much of the operational state that autonomous agents need to continue previous work (e.g., failed experiments, intermediate evidence, executable specifications, and implementation details). For RSI, such information loss directly weakens cross-generation knowledge accumulation.

To address this, they propose Agent-Native Research Artifact (ARA) as an executable alternative to conventional scientific documentation. ARA jointly preserves formal scientific logic, executable code and specifications, the exploration graph of both successful and abandoned branches, and raw empirical evidence. Rather than inheriting only a polished final narrative, a successor agent can therefore inspect how a result was obtained, which alternatives failed, and which assumptions or parameters were actually used. The reported evaluation shows 93.7\% question-answering accuracy over prior work using ARA compared with 72.4\% using conventional papers, while RE-Bench reproduction success increases from 57.4\% to 64.4\%. The important RSI contribution is thus not merely better documentation, but a richer inheritance substrate for later research cycles.

$\bullet$ \textbf{(2) Deterministic Verification for Research Claims.} Additionally, as autonomous agents generate experiments and claims at machine speed, \textit{verification rather than generation becomes the main bottleneck}. Its rit protocol anchors empirical claims directly to execution traces, re-extracting reported results from logs, while analytical claims can be checked through formal verification systems such as Lean 4. Only claims that pass these machine-verifiable gates are admitted into the shared research state. This prevents hallucinated or incorrectly reported results from being recursively inherited and amplified by later generations. They also augment sparse outcome-based optimization with epistemic-progress guidance and structured priors derived from human scientific reasoning. Instead of rewarding only the final benchmark score, agents are encouraged to prefer experiments that reduce uncertainty and reveal useful structure in the problem. This is intended to improve research efficiency in large hypothesis spaces, where naive hill climbing can repeatedly exploit familiar parameters without producing deeper understanding.

\textbf{Measured Improvement.} These ideas are instantiated in silicon design, where the feedback loop is unusually fast and deterministic. Autonomous agents generate synthesizable SystemVerilog, construct testbenches, and propose microarchitectural modifications; candidates are then evaluated using industrial EDA tools for synthesis, place-and-route, timing analysis, and formal equivalence. Invalid designs are filtered automatically, while verified implementations, execution traces, failed branches, and superior power-performance-area trajectories are retained as training and research evidence for subsequent generations. On Chip-Bench, they report a Level-3 CPU optimization score of 5.8416, with a design that is 2.7\% faster and 21.5\% smaller in area than standard-agent baselines on the same model foundation while passing all 97 directed tests and 350 random differential programs.

%As an RSI best practice, the case of Agent-Native Research Lab highlights two requirements that are often missing from self-improving systems: high-fidelity inheritance and trustworthy promotion. The research artifact preserves enough state for later agents to build on both successes and failures, while deterministic verification controls which claims are allowed to enter the lineage. The lab further maintains authenticated generation lineage, strict holdout separation, and a frozen baseline to isolate the effect of the evolving research artifact. Multi-generation compounding under matched budgets is still an ongoing evaluation target, but the system provides a concrete infrastructure for measuring whether autonomous research improvements genuinely accumulate rather than merely producing isolated better outputs.






\iffalse
Achieving genuine recursive self-improvement requires an AI-native infrastructure engineered from first principles for autonomous AI scientists. 
Today, every convention carrying scientific knowledge arose around human cognitive limits. A paper is linear because human attention operates serially. It relies on prose because natural language provides shared communication without shared tooling. It runs a few thousand words to match a referee's finite reading time. Peer review anchors trust in personal reputation because manual re-execution is prohibitively expensive. These historical conventions form an interface built specifically for human biology. Autonomous agents operating in a recursive self-improvement loop represent an entirely different species of researcher, proposing, executing, verifying, and inheriting findings at machine speed. Forcing these agents to communicate through human-tailored documents creates severe friction. Genuine recursive progress demands a scientific infrastructure built natively for artificial intelligence: redesigning how research is recorded, verified, and compounded across cycles, while preserving the foundational principles of scientific inquiry.

Human civilization began accumulating wisdom only after developing language. Language provided the transmission protocol that allowed insights to compound across generations, preserving discoveries across lifetimes. In autonomous research, the protocol used to document knowledge serves as this exact primitive for compounding. Traditional scientific papers break this accumulation mechanism by discarding the operational history an agent needs to build on past cycles. On RE-Bench, failed runs consume 90.2\% of dollar costs with a median failed-to-success token ratio of 113x, yet traditional write-ups drop these dead ends entirely. On PaperBench, only 45.4\% of 8,921 reproduction requirements are fully specified. When an inheritance loop drops nine-tenths of its experimental history and leaves half of its parameters implicit, knowledge compounding collapses. The Agent-Native Research Artifact (ARA) defines an executable transmission protocol for machine agents across four interlocking layers: formal scientific logic, code with executable specifications, an exploration graph preserving abandoned branches, and raw empirical evidence. On this protocol, agents achieve 93.7\% question-answering accuracy on prior work compared to 72.4\% on papers, and lift RE-Bench reproduction success from 57.4\% to 64.4\%.

Trust requires an equally fundamental overhaul because autonomous agents generate claims orders of magnitude faster than conventional science can evaluate them. When an agent runs unsupervised for days, generation becomes trivial while verification becomes the defining bottleneck. Language models are probabilistic and routinely hallucinate findings. Without rigorous verification, agents accept fabricated metrics as true, building successive experiments on invalid premises. These errors compound across generations, spinning the agents into endless cycles of ungrounded iteration that squander compute. Auditing these claims through prose or secondary model judges fails because evaluating one probabilistic model with another preserves the underlying uncertainty. Automated research demands deterministic verification. The \texttt{rit} protocol anchors claims directly to machine checks: empirical figures are deterministically re-extracted from execution logs, while analytical propositions pass through formal verification kernels like Lean 4. Claims enter the shared knowledge base only after passing these automated gates. Establishing verification as a strict gate ensures agents iterate on reproducible findings, preventing error cascades and protecting compute resources.

Much of the wisdom guiding scientific discovery can be learned directly from human inquiry and its underlying philosophy. Centuries of human science show that breakthroughs require more than brute-force optimization against a sparse outcome score. An agent optimizing solely for a final benchmark score repeatedly overfits familiar parameters. Productive exploration relies on epistemic progress, a concept rooted in the philosophy of science where researchers track how each intermediate experiment reduces uncertainty about the physical world. In simulated physics environments, agents equipped with a continuous epistemic-progress signal uncover hidden physical laws where traditional optimization baselines stall. Human inquiry also contributes the philosophy of research taste: the heuristic discernment that identifies promising hypotheses early. Frameworks like Sci-Reasoning decode the decision paths of skilled human scientists into structured priors that agents can inherit. Autonomous systems can adopt these human heuristics to guide their search through vast hypothesis spaces.

We applied these techniques to silicon design because hardware engineering provides an exceptionally fast, deterministic domain to close the self-improvement loop. Our work extends beyond standard model pretraining into an active recursive self-improvement flywheel. In this framework, autonomous agents write synthesizable SystemVerilog modules, generate rigorous testbenches, and implement microarchitectural optimizations. Every proposed implementation directly confronts industrial EDA toolchains: logic synthesis via Yosys, place-and-route under nextpnr, static timing analysis through OpenSTA, and formal equivalence proofs with SymbiYosys. Circuits must satisfy cycle-accurate correctness, timing closure, and strict physical area bounds. The toolchain filters out invalid attempts and automatically compiles physically verified solutions and superior PPA (power, performance, area) trajectories into high-quality training corpora. Training successive model generations on these verifiable engineering traces enables the system to produce progressively superior hardware implementations and search strategies over time.

We evaluate this self-evolving system on Chip-Bench (https://chipbench.ai/), an open benchmark maintained by ARA Labs for AI chip engineering. Chip-Bench evaluates agents across two demanding tracks governed entirely by EDA toolchains with zero model grading: a three-level CPU generation track culminating in full PPA optimization, and a hill-climb track containing 20 optimization tasks across 6 layers of the silicon stack. Standard agent models struggle against these physical constraints, frequently breaking cycle-level functional equivalence or failing timing closure during aggressive microarchitectural edits. In contrast, our self-evolving system, integrating continuous epistemic-progress guidance and structured exploration graphs, consistently outperforms standard agent models across the optimization suite. On Level 3 CPU optimization, the guided system achieves the highest public optimization score of 5.8416, delivering a processor core 2.7\% faster with 21.5\% less area compared to standard agent baselines on the same model foundation, while passing all 97 directed tests and 350 random differential programs. Across the hill-climb optimization suite, the system establishes leading records on challenging sub-tasks, including arithmetic pipeline frequency scaling (scoring 1.060 against reference) and logic synthesis recipes. These empirical gains demonstrate how verifiable, closed-loop evolution enables automated systems to solve dense engineering challenges where raw model generation falters.

The scope of these empirical results remains specific. Each produced artifact is preserved, guides subsequent work cycles, and supports independent verification without relying on the producer's credibility. Demonstrating matched-budget compounding across multiple generations remains an ongoing effort. The evaluation machinery is operational: authenticated lineage envelopes bind each generation to its parent state, benchmark reporting enforces strict holdout separation, and a frozen baseline isolates the research artifact as the single mutable variable. These instruments provide the foundation needed to measure multi-generation compounding curves systematically.
\fi
